\section{Introduction}
\label{sec:intro}

% 写作顺序建议：先写 03-taxonomy 和 04-evidence，最后回来写引言。
% 引言的核心任务是【提出问题并说明本文如何回答】，不是罗列文献。

\subsection{Motivation}

Large language models have become the substrate for a generation of agent
systems that plan, call tools, and act over extended horizons. Progress in
capability, however, has not come from parameters alone. Much of it comes from
the surrounding \emph{harness}: the scaffold that decides what context the model
sees, which tools it may call, how failures are handled, and when the loop
terminates \todo{cite survey 2508.07407 and 2607.13104 here}.

This raises a question that motivates the present review. If the harness
accounts for a substantial share of observed agent capability, and if the
harness is itself an ordinary program, then what happens when a system is
allowed to \emph{modify its own harness}? Does this constitute a form of
recursive self-improvement, or does it merely produce a better program once,
with no recurrence?

\subsection{Scope and contributions}

This review makes three contributions.

\begin{enumerate}
    \item \textbf{A five-axis taxonomy.} We organize \todo{N} primary works
    along axes that are independent of one another, chosen so that each paper
    can be assigned a position without reference to the other assignments.
    Section~\ref{sec:taxonomy} develops this framework.

    \item \textbf{A method comparison.} We compare the papers along nine
    dimensions, including two that the literature reports inconsistently:
    whether a compute-matched baseline was used, and whether search cost was
    accounted for. See Section~\ref{sec:evidence} and
    Table~\ref{tab:comparison}.

    \item \textbf{An identification of open problems.} We isolate
    \todo{2 or 3} questions that the current literature raises but does not
    resolve --- in particular, whether reported gains survive scrutiny for
    selection bias and whether the verifier is independent of the system being
    optimized. Section~\ref{sec:open}.
\end{enumerate}

\subsection{Exclusion criteria}

We exclude work that improves prompts without any element of recurrence
(prompt optimization in isolation), and work that modifies model parameters or
training data \todo{state how this boundary is drawn against the Model and
Data tracks}. We treat surveys as scaffolding for the domain map rather than as
primary evidence, and cite them accordingly.

\subsection{Organization}

Section~\ref{sec:background} fixes terminology. Section~\ref{sec:taxonomy}
presents the taxonomy. Section~\ref{sec:evidence} develops the evidence.
Section~\ref{sec:open} states the open problems.
Section~\ref{sec:conclusion} concludes.
